\begin{abstract}

碳化硅（SiC）外延层厚度的精确测量是功率半导体器件制造中的关键质量控制环节。本文基于红外反射光谱数据，围绕双光束干涉模型、三域互证融合算法及多光束干涉效应展开系统研究，实现了无标定参照条件下的 SiC 外延层厚度稳健估计。

对于问题一，建立了三层介质（空气/外延层/衬底）双光束干涉数学模型。从菲涅耳反射公式出发，导出了反射率谱的显式表达式 $R(\sigma)=r_{01}^2+\rho_{12}^2+2r_{01}\rho_{12}\cos(2\pi p\sigma+\varphi_{12})$ 及条纹级数与厚度的线性关系 $m=4n_1d\cos\theta_1\cdot\sigma_m+\varphi_{12}/\pi$。通过四项数据实验（EXP-1$\sim$4），以判据先行方式确定了模型参数：折射率常数 $n_1=2.55$、衬底界面有效相位 $\varphi_{12}\approx-0.28\pi$、反射幅值 $\rho_{12}\approx0.165$、标量偏振处理。

对于问题二，设计了三域互证融合算法体系：极值回归法（E1，光谱域，与 ASTM F95 同源）、谱域 FFT 法（E2，二次频域）、全谱拟合法（E3，模型驱动）。三算法分别在两个入射角（10$^\circ$、15$^\circ$）上独立估计后经逆方差加权融合，并通过两级一致性检验（L1 估计器间、L2 角度间）验证。分块 bootstrap（$N=200$）实证统计不确定度。结果表明，碳化硅外延层厚度为 8.03\,$\mu$m（统计不确定度 $\pm0.05\,\mu$m，模型系统区间 $[8.0,8.9]\,\mu$m），双入射角互证一致（相对差 0.4\%）。算法经合成数据闭环验证，恢复偏差 $\leq0.06\,\mu$m。

对于问题三，将双光束模型升级为 Airy 多光束模型，显式引入衬底 Drude 介电函数的频变相位。推导了可观测多光束干涉的四必要条件，经三重数值验证（闭式-逐束求和等价性、截断误差阶数 $O(q)$、倍频签名）确认模型正确性。对 Si 样品（附件3/4），三条证据链（残差达噪声 250 倍、条纹可见度衰减 24$\sim$28 倍、厚度差 51$\sigma$）一致判定出现显著多光束干涉，测得 Si 外延层厚度 $3.43\pm0.005_{\text{stat}}\pm0.15_{\text{sys}}\,\mu$m。对 SiC 样品，$2\times2$ 效应分离表明多光束可忽略（$q\approx0.002\sim0.005$），残差主因为衬底 Drude 相位频变；v2.0 修正 $+0.61\,\mu$m（$8.03\rightarrow8.64\,\mu$m），落在 v1.5 系统区间上缘。

本文的模型体系具有数学连续性（v2.0 在 $q\to0$ 时精确退化为 v1.0），验证体系覆盖算法闭环、数据互证、残差诊断、多算法交叉四个维度，方法学以判据先行、负结果入预算、可辨识性意识为特色。研究成果可为红外反射法测量外延层厚度提供模型选择依据与不确定度量化框架。

\keywords{红外反射；外延层厚度；双光束干涉；三域互证；多光束干涉；参数可辨识性；不确定度}

\end{abstract}